% Part: first-order-logic
% Chapter: models-theories
% Section: expressing-relations

\documentclass[../../../include/open-logic-section]{subfiles}

\begin{document}

\olfileid{fol}{mat}{exr}

\olsection{在结构中表达关系}

\begin{explain}
公式的一个主要用途是，使用结构~$\Struct M$ 的语言~$\Lang L$ 的初始符号，来表达~$\Struct M$ 中的性质与关系。
具体而言：$\Struct M$ 的论域是一个对象的集合。
常项符号、函数符号和谓词符号在~$\Struct M$ 中被分别解释为~$\Domain M$ 中的某些对象、$\Domain M$ 上的某些函数以及~$\Domain M$ 上的某些关系。
例如，如果 $\Lang L$ 中有 $\Obj A^2_0$，则 $\Struct M$ 会将某个关系~$R = \Assign{{\Obj
  A^2_0}}{M}$ 赋给它。
此时，公式 $\Atom{\Obj A^2_0}{\Obj v_1, \Obj v_2}$ 就\emph{表达}了该关系，其含义如下：若变元指派~$s$ 将 $\Obj v_1$ 映射到 $a \in \Domain{M}$，将 $\Obj v_2$ 映射到 $b \in \Domain M$，则
\[
Rab \text{\quad 当且仅当 \quad} \Sat{M}{\Atom{\Obj A^2_0}{\Obj v_1, \Obj v_2}}[s].
\]
注意，这里必须涉及变元指派：我们不能简单地说“$Rab$ 当且仅当 $\Sat{M}{\Atom{\Obj A^2_0}{a, b}}$”，因为 $a$ 和 $b$ 不是我们语言中的符号，它们是~$\Domain{M}$ 中的元素。

由于我们不仅有原子公式，还可以用逻辑联结词和量词组合它们，因此更复杂的公式可以定义那些并非~$\Struct M$ 中直接给定的其他关系。
我们关心的是如何做到这一点，特别地，我们能在结构中定义哪些关系。
\end{explain}

\begin{defn}
设 $!A(\Obj v_1,\dots, \Obj v_n)$ 是 $\Lang L$ 的一个公式，其中只有 $\Obj v_1$，\dots，$\Obj v_n$ 自由出现，并设 $\Struct M$ 是 $\Lang L$ 的一个结构。$!A(\Obj v_1,\dots, \Obj v_n)$
\emph{表达了关系}~$R \subseteq \Domain M^n$，当且仅当对任意满足 $s(\Obj v_i) = a_i$（$i = 1,
\dots, n$）的变元指派~$s$，都有
\[
Ra_1\dots a_n \text{\quad 当且仅当 \quad} \Sat{M}{\Atom{!A}{\Obj
    v_1,\dots, \Obj v_n}}[s]
\]
\end{defn}

\begin{ex}
在算术的标准模型~$\Struct N$ 中，公式 $\Obj
v_1 < \Obj v_2 \lor \eq[\Obj v_1][\Obj v_2]$ 表达了 $\Nat$ 上的 $\le$ 关系。
公式 $\eq[\Obj v_2][\Obj v_1']$ 表达了后继关系，即关系 $R \subseteq \Nat^2$，其中 $Rnm$ 成立当且仅当 $m$ 是 $n$ 的后继。
公式 $\eq[\Obj v_1][\Obj v_2']$ 表达了前驱关系。
公式 $\lexists[\Obj v_3][(\eq/[\Obj v_3][\Obj 0] \land
  \eq[\Obj v_2][(\Obj v_1 + \Obj v_3)])]$ 和 $\lexists[\Obj
  v_3][\eq[(\Obj v_1 + {\Obj v_3}')][v_2]]$ 都表达了 $\Obj <$ 关系。
这意味着在算术语言中，谓词符号~$<$ 实际上是多余的；它是可定义的。
\end{ex}

\begin{explain}
这一思想不仅在特定结构中有意义，在一般情形下——即每当我们使用语言来描述一个或多个预期的模型（亦即考虑理论）时——也同样有意义。
这些理论通常只包含少数几个谓词符号作为基本符号，但在它们所描述的论域中，常常有许多其他关系扮演着重要角色。
如果这些其他关系能够被那些解释了语言中基本谓词符号的关系系统地表达出来，我们就说可以在该语言中\emph{定义}它们。
\end{explain}

\begin{prob}
在 $\Lang L_A$ 中找出定义下列关系的公式：
\begin{enumerate}
\item $n$ 在 $i$ 和 $j$ 之间；
\item $n$ 整除 $m$（即 $m$ 是 $n$ 的倍数）；
\item $n$ 是素数（即除了 $1$ 和 $n$ 之外没有其他数能整除~$n$）。
\end{enumerate}
\end{prob}

\begin{prob}
假设公式 $!A(\Obj v_1, \Obj v_2)$ 在结构~$\Struct M$ 中表达了关系 $R \subseteq \Domain M^2$。请找出表达下列关系的公式：
\begin{enumerate}
\item $R$ 的逆 $R^{-1}$；
\item $R$ 的相对积 $R \mid R$；
\end{enumerate}
你能找到表达 $R$ 的传递闭包 $R^+$ 的方法吗？
\end{prob}

\begin{prob}
令 $\Lang{L}$ 为仅包含一个二元谓词符号 $<$ 的语言（没有其他常项符号、函数符号或谓词符号——当然，等词~$\eq$ 除外）。
令 $\Struct{N}$ 为满足 $\Domain{N} = \Nat$ 且 $\Assign{<}{N} = \Setabs{\tuple{n,m}}{n <
  m}$ 的结构。证明下列命题：
\begin{enumerate}
\item $\{ 0 \}$ 在 $\Struct{N}$ 中可定义；
\item $\{ 1 \}$ 在 $\Struct{N}$ 中可定义；
\item $\{ 2 \}$ 在 $\Struct{N}$ 中可定义；
\item 对每个 $n \in \Nat$，集合 $\{ n \}$ 在 $\Struct{N}$ 中可定义；
\item $\Domain{N}$ 的每个有限子集在 $\Struct{N}$ 中可定义；
\item $\Domain{N}$ 的每个余有限子集在 $\Struct{N}$ 中可定义（其中 $X \subseteq \Nat$ 是余有限的当且仅当 $\Nat \setminus X$ 是有限的）。
\end{enumerate}
\end{prob}

\end{document}
